% THRES-LOSSYPROBE2 ADDENDUM 2026-10-18
\documentclass[11pt]{article}
\usepackage{amsmath,amssymb,amsthm}
\newtheorem{theorem}{Theorem}
\title{THRES-LOSSYPROBE2: Guarded Restriction and Direct-Sum Controls}
\date{18 October 2026}
\begin{document}\maketitle
\noindent\textbf{Status.} PROVED/EXACT in declared models,
pending independent review. No unconditional P/NP separation.
\section{Guarded hard subcubes}
For $N=2^n$ and circuit threshold $s$ define
$B(m,s)$ as an upper bound on the base-2 logarithm
of the number of $m$-input circuits with at most $s$ gates:
\[
 B(m,s)=\lceil\log_2(s+1)\rceil+
 \lceil\log_2(m+s)\rceil+
 s(2+2\lceil\log_2(m+s)\rceil).
\]
\begin{theorem}[SLICE-HIDE2]
Fix $k$ designated guard variables, $m=n-k$,
and assume $2^m>B(m,s)$. For any threshold $r<s$,
an exact algorithm whose only observations are
MCSP[$r$] queries on restrictions fixing all $k$
guard variables must query all $2^k$ guard patterns
to decide parent MCSP[$s$] on every truth table.
The same holds if each query returns the entire subtable.
\end{theorem}
\noindent Proof. Counting yields $h$ on $m$ inputs
with $\operatorname{CC}(h)>s$. On the zero parent,
let $a_\star$ be a guard assignment not queried.
Define $F(u,v)=[u=a_\star]\wedge h(v)$.
Setting $u=a_\star$ shows
$\operatorname{CC}(F)\ge\operatorname{CC}(h)>s$.
All queried restrictions of $F$ are identically zero,
so the adaptive transcript is identical to that on
the zero parent. Decisions must differ, contradiction.
For $s=N^\beta$, fixed $\beta<1$, the choice
$k=\lfloor(1-\beta)n/2\rfloor$ satisfies the
assumption for sufficiently large $n$.
For a uniformly hidden guard assignment,
$q$ adaptive same-guard queries hit the special
pattern with probability at most $q/2^k$.
Full-source access or arbitrary other restrictions
escape the theorem and must be charged independently.
\section{Pairwise distant yet globally shared}
\begin{theorem}[DISTANT-BLOCK-SHARE2]
For arbitrary $f:\{0,1\}^n\to\{0,1\}$,
let $q=2^n$ and
$f_i(x)=f(x)\oplus\langle i,x\rangle_{\mathbb F_2}$
for all $i\in\{0,1\}^n$.
For any distinct $i,j$,
$\operatorname{dist}_H(f_i,f_j)=2^{n-1}$.
Furthermore,
\[
 \operatorname{CC}(f_i)\ge\operatorname{CC}(f)-4n,
 \qquad
 \operatorname{CC}(F)\le\operatorname{CC}(f)+5n,
 \quad F(i,x)=f_i(x).
\]
Hence all $q$ blocks may be individually hard and
pairwise Hamming distance $1/2$ while the joint
indexed function shares one hard core at $O(n)$ overhead.
\end{theorem}
\noindent Proof. The difference $f_i\oplus f_j$
is the balanced nonzero linear form
$\langle i\oplus j,x\rangle$.
The identity
$a\oplus b=(a\vee b)\wedge\neg(a\wedge b)$
uses four AND/OR/NOT gates.
For fixed $i$, $f$ is recovered from $f_i$
with at most $4n$ more gates.
For $F$, compute the $n$ products $i_j\wedge x_j$,
then their parity and one final XOR with $f(x)$,
using $n+4(n-1)+4=5n$ additional gates.
Take any $f$ with $\operatorname{CC}(f)>N^\beta$,
which exists for fixed $\beta<1$ and large $n$.
Then the sum of the $q$ individual circuit sizes
is $\Omega(q\operatorname{CC}(f))$ but
$\operatorname{CC}(F)=O(\operatorname{CC}(f)+n)$.
Pairwise distinctness, high complexity, and
exact half-distance are insufficient to prove
a generic direct-sum complexity lower bound.
\section{Proof-track consequence}
The remaining THRESHOLD-DESCENT1 arrow needs
a lossy YES/NO-preserving quotient with full
input-sweep accounting and a same-one-tape
target-to-source simulation. Pairwise block
distance and oblivious guard restrictions do
not supply it. An index-hiding or conditional
information theorem must defeat this shared-core family.
\end{document}
% END_OF_VERIFIED_REPORT::THRES-LOSSYPROBE2-ADDENDUM::P_VS_NP_OPEN
